\chapter{Generaci\'{o}n de c\'{o}digo}
\label{cap:conclusiones}

La generaci\'{o}n de c\'{o}digo se realiza a trav\'{e}s del main de constructorast pasandole el nombre del fichero como par\'{a}metro
y subiendo el archivo .wat generado a la carpeta generador. Despu\'{e}s convertimos el fichero .wat a wasm y lo ejecutamos con la funci\'{o}n de javascript dada.




Los ejemplos de generaci\'{o}n de c\'{o}digo est\'{a}n divididos. Por una parte, el c\'{o}digo fuente se encuentra en la carpeta de test, donde empiezan a partir de input20.txt ya que el resto son test de detecci\'{o}n de errores.
Por otra parte, el c\'{o}digo .wat generado se va a la carpeta de generador donde tambi\'{e}n se encuentra el main.js dado en ejercicios.

\section{Ejemplos de c\'{o}digo}

\subsection{Ejemplo 1}
Input20.txt en la carpeta de tests, es un ejemplo sencillo el cual puede compilarse y posteriormente probarse en la carpeta de generador.

\begin{algorithm}

	\Integer{} x  \;
	\Integer{} y  \;
	
	\Void() \reserved{main}() \{ \\
	\quad x := \const{0} \;
	\quad y := y + \const{1}\;
	\quad \reserved{prints}(y)\;
	\}\;

 	\Integer{} z \;
	\Integer{} w \;
\end{algorithm}



\subsection{Ejemplo 2}
Input21.txt en la carpeta de tests, es un ejemplo sencillo parecido al anterior pero se agrega una llamada a funci\'{o}n.
\begin{algorithm}

	\Integer{} x  \;
	\Integer{} y  \;
	
	\Void() \reserved{main}() \{ \\
	\quad x := \const{0} \;
	\quad y := y + \const{1}\;
	\quad x := f(y)\;
	\quad \reserved{prints}(y)\;
	\quad \reserved{prints}(x)\;
	\}\;

 	\Integer{} z \;
	\Integer{} w \;
	
	\reserved{int}{(\reserved{int})} f(y) \{ \\
	\quad z := y \;
	\quad \reserved{prints}(z) \;
	\quad y := y + \const{1} \;
	\quad \reserved{prints}(y) \;
	\quad \reserved{return} \const{2}*z \;
	\}\;
	
\end{algorithm}

\newpage

\subsection{Ejemplo 3}
Input22.txt en la carpeta de tests, es un ejemplo sencillo parecido al anterior pero se agrega una variable local.

void() main(){
    x := 0;
    y := x + 1;
    int r;
    scans(r);
    prints(r);
    x := f(r);
    prints(y);
    prints(r);
    prints(y);
    prints(x);
};
int z;
int w;
int(int) f(y){ //oculta a la global y
    z := y;
    prints(z);
    y := y + 3;
    prints(y);
    return 2*z;
};
\begin{algorithm}

	\Integer{} x  \;
	\Integer{} y  \;
	
	\Void() \reserved{main}() \{ \\
	\quad x := \const{0} \;
	\quad y := y + \const{1}\;
	\quad int r\;
	\quad x := f(y)\;
	\quad \reserved{prints}(r)\;
	\quad \reserved{prints}(y)\;
	\quad \reserved{prints}(x)\;
	\}\;

 	\Integer{} z \;
	\Integer{} w \;
	
	\reserved{int}{(\reserved{int})} f(y) \{ \\
	\quad z := y \;
	\quad \reserved{prints}(z) \;
	\quad y := y + \const{3} \;
	\quad \reserved{prints}(y) \;
	\quad \reserved{return} \const{2}*z \;
	\}\;
	
\end{algorithm}

\newpage

\subsection{Ejemplo 4}
Input23.txt en la carpeta de tests, es un ejemplo sencillo parecido al anterior pero se lleva por valor la variable global.

\begin{algorithm}

	\Integer{} x  \;
	\Integer{} y  \;
	
	\Void() \reserved{main}() \{ \\
	\quad x := \const{0} \;
	\quad y := y + \const{1}\;
	\quad int r\;
	\quad x := f(y)\;
	\quad \reserved{prints}(r)\;
	\quad \reserved{prints}(y)\;
	\quad \reserved{prints}(x)\;
	\}\;

 	\Integer{} z \;
	\Integer{} w \;
	
	\reserved{int}{(\reserved{int})} f(y) \{ \\
	\quad z := y \;
	\quad \reserved{prints}(z) \;
	\quad y := y + \const{3} \;
	\quad \reserved{prints}(y) \;
	\quad \reserved{return} \const{2}*z \;
	\}\;
	
\end{algorithm}

\newpage


\subsection{Ejemplo 5}
Input24.txt en la carpeta de tests, es un ejemplo sencillo parecido al anterior pero se usan arrays.

\begin{algorithm}

	\Integer{} x  \;
	\Integer{} y  \;
	
	\Void() \reserved{main}() \{ \\
	\quad x := \const{0} \;
	\quad y := y + \const{1}\;
	\quad int n\;
	\quad \reserved{scans}(n)\;
	\quad \reserved{prints}(n)\;
	\quad \reserved{int}[\const{3}] (a)\;
	\quad a[\const{0}] := x\;
	\quad a[\const{1}] := y\;
	\quad a[\const{2}] := y + \const{1}\;
	\quad x := f(a[n])\;
	\quad \reserved{prints}(a[\const{2}])\;
	\quad \reserved{prints}(a[n])\;
	\}\;

 	\Integer{} z \;
	\Integer{} w \;
	
	\reserved{int}{(\reserved{int})} f(y) \{ \\
	\quad z := y \;
	\quad \reserved{prints}(z) \;
	\quad y := y + \const{3} \;
	\quad \reserved{prints}(y) \;
	\quad \reserved{return} \const{2}*z \;
	\}\;
	
\end{algorithm}

\newpage

\subsection{Ejemplo 6}
Input25.txt en la carpeta de tests, es un ejemplo sencillo parecido al anterior pero se usan structs.

\begin{algorithm}

	\Integer{} x  \;
	\Integer{} y  \;
	
	\reserved{struct} A \{
		\quad \reserved{int} a\,
		\quad \reserved{int} b\;
	\} \;
	
	\Void() \reserved{main}() \{ \\
	\quad x := \const{0} \;
	\quad y := y + \const{1}\;
	\quad int n\;
	\quad \reserved{scans}(n)\;
	\quad \reserved{prints}(n)\;
	\quad \reserved{struct} A s\;
	\quad s.b := n\;
	\quad s.a := 2\;
	\quad \reserved{prints}(s.a)\;
	\quad x := f(s.a)\;
	\quad \reserved{prints}(s.b)\;
	\quad \reserved{prints}(s.a)\;
	\}\;

 	\Integer{} z \;
	\Integer{} w \;
	
	\reserved{int}{(\reserved{int})} f(y) \{ \\
	\quad z := y \;
	\quad \reserved{prints}(z) \;
	\quad y := y + \const{3} \;
	\quad \reserved{prints}(y) \;
	\quad \reserved{return} \const{2}*z \;
	\}\;
	
\end{algorithm}

\newpage


\subsection{Ejemplo 7}
Input26.txt en la carpeta de tests, es un ejemplo que combina los dos anteriores.

\begin{algorithm}

	\Integer{} x  \;
	\Integer{} y  \;
	
	\reserved{struct} A \{
		\quad \reserved{int}[\const{3}] a\,
		\quad \reserved{int} b\;
	\} \;
	
	\Void() \reserved{main}() \{ \\
	\quad x := \const{0} \;
	\quad y := y + \const{1}\;
	\quad int n\;
	\quad \reserved{scans}(n)\;
	\quad \reserved{prints}(n)\;
	\quad \reserved{struct} A s\;
	\quad s.b := n\;
	\quad s.a[\const{0}] := x\;
	\quad s.a[\const{1}] := y\;
	\quad s.a[\const{2}] := y + \const{1}\;
	\quad x := f(a[n])\;
	\quad \reserved{prints}(s.a[\const{n}])\;
	\quad \reserved{prints}(s.b)\;
	\quad \reserved{prints}(s.a[n])\;
	\}\;

 	\Integer{} z \;
	\Integer{} w \;
	
	\reserved{int}{(\reserved{int})} f(y) \{ \\
	\quad z := y \;
	\quad \reserved{prints}(z) \;
	\quad y := y + \const{3} \;
	\quad \reserved{prints}(y) \;
	\quad \reserved{return} \const{2}*z \;
	\}\;
	
\end{algorithm}










